\documentclass[12pt]{article}

\usepackage[utf8]{inputenc}
\usepackage{amsmath, amssymb, amsthm}
\usepackage{geometry}
\geometry{a4paper, margin=2.5cm}

\title{Teorema Banach--Alaoglu}
\author{}
\date{}

\theoremstyle{definition}
\newtheorem{definition}{Definiție}

\theoremstyle{plain}
\newtheorem{theorem}{Teoremă}
\newtheorem{proposition}{Propoziție}

\begin{document}

\maketitle

\begin{definition}
Fie $X$ un spațiu normat. Notăm prin $X^\ast$ spațiul dual, adică mulțimea tuturor funcționalelor liniare continue $f : X \to \mathbb{R}$ (sau $\mathbb{C}$), înzestrat cu norma


\[
\|f\| = \sup_{\|x\|\le 1} |f(x)|.
\]


Pentru $M > 0$, mulțimea


\[
B_{X^\ast}(0,M) = \{ f \in X^\ast : \|f\| \le M \}
\]


se numește \emph{discul închis de rază $M$ din dual}.
\end{definition}

\begin{theorem}[Banach--Alaoglu]
Fie $X$ un spațiu normat. Atunci discul unitar închis din dual,


\[
B_{X^\ast}(0,1) = \{ f \in X^\ast : \|f\| \le 1 \},
\]


este compact în topologia $\sigma(X^\ast, X)$, adică în topologia slab‑$\ast$ generată de aplicațiile


\[
f \mapsto f(x), \qquad x \in X.
\]


\end{theorem}

\begin{proposition}
În ipotezele teoremei Banach--Alaoglu, următoarele afirmații sunt echivalente pentru o mulțime $K \subset X^\ast$:
\begin{enumerate}
    \item[(i)] $K$ este slab‑$\ast$ compact.
    \item[(ii)] $K$ este închisă în topologia slab‑$\ast$ și este conținută într-un disc închis
    

\[
    K \subset B_{X^\ast}(0,M)
    \]


    pentru un anumit $M > 0$.
\end{enumerate}
\end{proposition}

\begin{proposition}[Formă secvențială a teoremei Banach--Alaoglu]
Fie $X$ un spațiu normat (separabil) și fie $(f_n)_{n\in\mathbb{N}}$ o familie de funcționale liniare continue pe $X$ astfel încât există $M>0$ cu


\[
\|f_n\|\le M \qquad \text{pentru toate } n\in\mathbb{N}.
\]


Atunci există un subsir $(f_{n_k})_{k\in\mathbb{N}}$ și un funcțional $f\in X^\ast$ cu $\|f\|\le M$ astfel încât


\[
f_{n_k}(x)\longrightarrow f(x) \qquad \text{pentru orice } x\in X.
\]


Cu alte cuvinte, orice șir mărginit în $X^\ast$ admite un subsir convergent în topologia slab‑$\ast$.
\end{proposition}


\end{document}
